\section{Schubert coefficients of disconnected matroids}

In this section we prove the following result, expressing the Schubert coefficients of a disconnected matroid in terms of the Schubert coefficients of its connected components.
\begin{thm}
    \label{thm:disconnected}
    Let $M_i$ be a matroid of rank $r_i$ on $[n_i]$ for $i = 1, 2$. The class of $M_1\oplus M_2$ in $A^\bullet(G(r_1+r_2,n_1+n_2))$ is
    \begin{equation*}
        \Sc(M_1\oplus M_2) = (\square \Sc(M_1))(\square \Sc(M_2))
    \end{equation*}
    where
    \begin{equation*}
        \square \Sc(M_i) = \sum_\mu d_\mu(M_i) \sigma_{\square \mu},
    \end{equation*}
    and \(\square \mu\) is obtained from \(\mu\) by attaching a \(r_i\) by \(n_{3-i} - r_{3-i}\) rectangle to the left of the Young diagram of \(\mu\).
\end{thm}
Notice that there is no assumption on \(M_i\) being connected in \cref{thm:disconnected}, so if \(M\) has more than \(2\) connected components applying the theorem repeatedly yields a formula for \(\Sc(M)\) in terms of the Schubert coefficients of its connected components.

\begin{example}
    We use \cref{thm:disconnected} to compute \(\Sc(U_{2,4} \oplus U_{2,5})\). By \cref{thm:Klyachko} we have
    \begin{align*}
        \Sc(U_{2,4}) &= 2 \sigma_{(1,0)}, \\
        \Sc(U_{2,5}) &= 3 \sigma_{(2,0)} + 1 \sigma_{(1,1)}.
    \end{align*}
    Attaching a \(2\times 3\) rectangle to the left of \((1,0)\) and a \(2 \times 2\) rectangle to the left of \((2,0)\) and \((1,1)\) gives
    \begin{align*}
        \square\Sc(U_{2,4}) &= 2 \sigma_{(4,3)}, \\
        \square \Sc(U_{2,5}) &= 3 \sigma_{(4,2)} + 1 \sigma_{(3,3)}.
    \end{align*}
    Now computing the product \((\square\Sc(U_{2,4}))(\square \Sc(U_{2,5}))\) in \(A^\bullet(G(4,9))\) gives
    \begin{align*}
        \Sc(U_{2,4} \oplus U_{2,5}) &= (2 \sigma_{(4,3)})(3 \sigma_{(4,2)} + 1 \sigma_{(3,3)})\\
        &= 2\sigma_{(4, 3, 3, 3)} + 8\sigma_{(4, 4, 3, 2)} + 6\sigma_{(4, 4, 4, 1)} + 8\sigma_{(5, 3, 3, 2)} \\
        & + 8\sigma_{(5, 4, 2, 2)} + 14\sigma_{(5, 4, 3, 1)} + 6\sigma_{(5, 4, 4, 0)} + 8\sigma_{(5, 5, 2, 1)} + 6\sigma_{(5, 5, 3, 0)}.
    \end{align*}
\end{example}

Throughout fix $r_1 < n_1$ and $r_2 < n_2$ and let $r = r_1 + r_2$ and $n = n_1 + n_2$. We consider the embedding 
\begin{equation*}
    \iota \colon G(r_1,n_1) \times G(r_2,n_2) \hookrightarrow G(r,n)
\end{equation*}
defined by $L(\iota(x,y)) = L(x) \oplus L(y) \subseteq \CC^{n_1} \oplus \CC^{n_2} = \CC^n$.

\begin{lemma}
    \label{lem:varietyLift}
    Let $X$ be a subvariety of $G(r_1,n_2)$ of codimension $s$ with
    \begin{equation*}
        [X] = \sum_\mu d_\mu(X) \sigma_\mu \in A^\bullet(G(r_1,n_1))
    \end{equation*}
    for $d_\mu(X) \in \ZZ$ and let $\tilde{X} = \iota(X \times G(r_2,n_2))$. Then
    \begin{equation*}
        [\tilde{X}] = \sum_\mu d_\mu(X) \sigma_{\square \mu} \in A^\bullet(G(r,n))
    \end{equation*}
    where both sums are over all partitions $\mu \subseteq r_1 \times (n_1-r_1)$ with $|\mu| = s$ and $\square \mu = (n_2-r_2 + \mu_1, \dots, n_2-r_2 + \mu_{k_1}) \subseteq r\times (n-r)$.
\end{lemma}
\begin{proof}
    For fixed $\mu$ we need to show that 
    \[
        \deg([\tilde{X}] \sigma_{(\square \mu)^c}) = \deg([X] \sigma_{\mu^c}) = d_\mu(X).
    \]
    Let $V_\bullet$ be a complete flag of $\CC^{n_1}$ such that the intersection $X \cap X_{\mu^c}(V_\bullet)$ is transverse. Fix a $k_2$-dimensional subspace $U$ of $\CC^{n_2}$ and let $V^+_\bullet$ be a complete flag of $\CC^n = \CC^{n_1} \oplus \CC^{n_2}$ such that $V^+_{k_2} = 0 \oplus U$ and  $V^+_{n_2 + j} = V_j \oplus \CC^{n_2}$ for $j = 1, 2, \dots, n_1$. Let $(j_i)_{i=1}^{r_1}$ be the jumping sequence of $\mu^c$ in $r_1 \times (n_1-r_1)$. Then notice that the complement of $\square \mu$ in $r \times (n-r)$ is the partition obtained by attaching the $r_2 \times (n-r)$ rectangle to the top of the Young diagram of $\mu^c$ i.e. the jumping sequence of $(\square \mu)^c$ is
    \begin{equation*}
        (1,2,\dots, r_2, n_2 + j_1, n_2 + j_2, \dots n_2 + j_{r_1}).
    \end{equation*}
    From the first part of the jumping sequence above only the last term contributes with a defining inequality of the corresponding Schubert variety. So $X_{(\square \mu)^c}(V^+_\bullet)$ is given by 
    \begin{equation*}
        \left\{z \in G(r,n) \mid \dim(L(z) \cap V^+_{r_2}) \geq r_2, \ \dim(L(z) \cap V^+_{n_2+j_i}) \geq r_2 + i\right\}.
    \end{equation*}
    We need to show that $d_\mu(X) = |X \cap X_{\mu^c}(V_\bullet)| = |\tilde{X} \cap X_{(\square \mu)^c}(V^+_\bullet)|$. Suppose $z \in \tilde{X} \cap X_{(\square \mu)^c}(V^+_\bullet)$, since $z \in \tilde{X}$,  $L(z) = L(x) \oplus L(y)$ for $x \in X$ and $y \in G(r_2,n_2)$. Now
    \begin{equation*}
        L(z) \cap V^+_{r_2} = (L(x) \oplus L(y)) \cap (0 \oplus U) = 0 \oplus (L(y) \cap U)
    \end{equation*}
    so the first defining inequality of $X_{(\square \mu)^c}(V^+_\bullet)$ becomes 
    \begin{equation*}
        \dim(0\oplus (L(y) \cap U)) \geq r_2 \iff L(y) = U
    \end{equation*}
    since both $L(y)$ and $U$ are of dimension $r_2$. For the second type of defining inequality of $X_{(\square \mu)^c}(V^+_\bullet)$ we have
    \begin{equation*}
        L(z) \cap V^+_{n_2 + j_i} = (L(x) \oplus L(y)) \cap (V_{j_i} \oplus \CC^{n_2}) = (L(x) \cap V_{j_i}) \oplus U.
    \end{equation*}
    so 
    \begin{equation*}
        \dim(L(z) \cap V^+_{n_2 + j_i}) \geq k_2 + i \iff \dim(L(x) \cap V_{j_i}) \geq i
    \end{equation*}
    which are the defining inequalities of $X_{\mu^c}(V_\bullet)$. Hence if $z = \iota(x,y) \in \tilde{X} \cap X_{(\square \mu)^c}(V^+_\bullet)$ then $x \in X \cap X_{\mu^c}(V_\bullet)$ and if $x \in X \cap X_{\mu^c}(V_\bullet)$ then $\iota(x,U) \in \tilde{X} \cap X_{(\square \mu)^c}(V^+_\bullet)$ so $|X \cap X_{\mu^c}(V_\bullet)| = |\tilde{X} \cap X_{(\square \mu)^c}(V^+_\bullet)|$.
\end{proof}

By a similar argument for subvarieties $Y$ of $G(r_2,n_2)$ with 
\begin{equation*}
    [Y] = \sum_\eta d_\eta(Y) \sigma_\eta \in A^\bullet(G(r_2,n_2)),
\end{equation*}
let $\tilde{Y} = \iota(G(r_1,n_2) \times Y)$. Then we have
\begin{equation*}
    [\tilde{Y}] = \sum_\eta d_\eta(Y) \sigma_{\square \eta} \in A^\bullet (G(r,n))
\end{equation*}
where $\square \eta = (n_1-r_1+\eta_1, \dots, n_1-r_1 + \eta_{r_2}) \subseteq r \times (n-r)$.

Using \cref{lem:varietyLift} we can now prove \cref{thm:disconnected} in the case where $M_1$ and $M_2$ are representable over $\CC$.
\begin{proof}
    Fix $x \in G(r_1,n_1)$ and $y \in G(r_2,n_2)$ and let $z = \iota(x,y)$. Let $T_i = (\CC^*)^{n_i}$ be the algebraic torus acting on $G(r_i,n_i)$  for $i = 1,2$ and $T = T_1 \times T_2 = (\CC^*)^n$ be the algebraic torus acting on $G(r,n)$. Notice that
    \begin{equation*}
        \overline{Tz} = \iota(\overline{T_1 x} \times \overline{T_2 y}) = \iota(\overline{T_1 x} \times G(r_2,n_2)) \cap \iota(G(r_1,n_2) \times \overline{T_2 y}),
    \end{equation*}
    hence $[\overline{Tz}] = [\tilde{\overline{T_1 x}}][\tilde{\overline{T_2 y}}]$.
    
    Since $M_z = M_x \oplus M_y$ applying \cref{lem:varietyLift} to $\overline{T_1 x}$ and $\overline{T_2 y}$ gives
    \begin{equation*}
        \Sc(M_z) = (\sum_\mu d_\mu(M_x) \sigma_{\square \mu}) (\sum_\eta d_\eta(M_y) \sigma_{\square \eta}) = (\square Sc(M_x))(\square Sc(M_y)). \qedhere
    \end{equation*}
\end{proof}

To prove \cref{thm:disconnected} for arbitrary matroids, we use the fact that valuative matroid invariants are determined by their values on Schubert matroids, which are representable over $\CC$. In particular if two valuative matroid invariants are equal on all matroids representable over $\CC$, they are equal on all matroids.

\begin{proof}[proof of \cref{thm:disconnected}]
    Fix a matroid $M_1 \in S(r_1,n_1)$ that is representable over $\CC$. Notice that both
    \begin{equation*}
        \Sc(M_1 \oplus - ) \colon S(r_2,n_2) \to A^\bullet(G(r,n))
    \end{equation*}
    and 
    \begin{equation*}
        (\square Sc(M_1)) (\square Sc(-)) \colon S(r_2,n_2) \to A^\bullet(G(r,n))
    \end{equation*}
    are valuative matroid invariants that are equal on matroids representable over $\CC$ hence $\Sc(M_1 \oplus M_2) = (\square \Sc(M_1)) (\square \Sc(M_2))$ for all matroids $M_2 \in S(r_2,n_2)$.
    
    Now fix a matroid $M_2 \in S(r_2,n_2)$ and notice again that 
    \begin{equation*}
        \Sc(- \oplus M_2) \colon S(r_1,n_1) \to A^\bullet(G(r,n))
    \end{equation*}
    and 
    \begin{equation*}
        (\square \Sc(-)) (\square \Sc(M_2)) \colon S(r_1,n_1) \to A^\bullet(G(r,n))
    \end{equation*}
    are valuative matroid invariants that are equal on matroids representable over $\CC$ hence $\Sc(M_1 \oplus M_2) = (\square \Sc(M_1)) (\square \Sc(M_2))$ for all matroids $M_1 \in S(r_1,n_1)$ and $M_2 \in S(r_2,n_2)$.
\end{proof}

\begin{corollary}
\label{cor:disconnected}
    Fix a partition $\lambda \subset r \times (n-r)$ with $|\lambda| = r(n-r) - (n - \kappa(M_1 \oplus M_2))$. Then the Schubert coefficient $d_\lambda(M_1 \oplus M_2)$ is given by
    \begin{equation*}
        d_\lambda(M_1 \oplus M_2) = \sum_{\mu, \eta} d_\mu(M_1) d_\eta(M_2) c_{\square\mu \square\eta}^\lambda
    \end{equation*}
    where $c_{\square\mu \square\eta}^\lambda$ is a Littlewood-Richardson coefficient.
\end{corollary}
\begin{proof}
    By \cref{thm:disconnected} we have
    \begin{align*}
        \Sc(M_1 \oplus M_2) &= (\sum_\mu d_\mu(M_1) \sigma_{\square \mu}) (\sum_\eta d_\eta(M_2) \sigma_{\square \eta})\\
        &= \sum_{\mu, \eta} d_\mu(M_1) d_\eta(M_2) \sigma_{\square \mu} \sigma_{\square \eta}
    \end{align*}
    and since $d_\lambda(M_1\oplus M_2) = \deg(\Sc(M_1 \oplus M_2) \sigma_{\lambda^c})$ we have
    \begin{align*}
        d_\lambda(M_1 \oplus M_2) &= \sum_{\mu, \eta} d_\mu(M_1) d_\eta(M_2) \deg(\sigma_{\square \mu} \sigma_{\square \eta} \sigma_{\lambda^c})\\
        &= \sum_{\mu, \eta} d_\mu(M_1) d_\eta(M_2) c_{\square\mu \square\eta}^\lambda
    \end{align*}
    All the sums above are over partitions $\mu \subseteq r_1 \times (n_1 - r_1)$ with $|\mu| = r_1(n_1-r_1) - (n_1 - \kappa(M_1))$ and $\eta \subseteq r_2 \times (n_2 - r_2)$ with $|\eta| = r_2(n_2-r_2) - (n_2 - \kappa(M_2))$.
\end{proof}

Notice that the formula for \(d_\lambda(M_1\oplus M_2)\) in \cref{cor:disconnected} preserves non-negativity of the Schubert coefficients. That is if both \(M_1\) and \(M_2\) have non-negative Schubert coefficients, so does \(M_1 \oplus M_2\). So it is enough to prove \cref{conj:positivity} for connected matroids.